\scanpage{7}
% archive foliation: 6
\unsure{Soit}
\[
(13)\quad \underset{\substack{\cup\\\pi_1}}{\mathcal E'}=\operatorname{Norm}_{\mathcal E}(N)
\]
$=$ stabilisateur dans $\mathcal E$ de la relation d'équivalence dans $\widetilde{\mathcal X}$ définie par $\widetilde{\mathcal X}\mathbin{-}\mathcal X'$ i.e. par op. de $N\ldots$
\[
(14)\qquad \mathcal G'=\mathcal E'/N
\]
de sorte qu'on a
\[
(15)\quad \mathcal G'\xrightarrow{\sim}\operatorname{Aut}_k(\mathcal X'\longrightarrow\mathcal X)
\]
et une suite exacte
\[
(16)\quad\diagram{\begin{tikzcd}[ampersand replacement=\&]
1\ar[r,no head]\&G\ar[r]\&\mathcal G'\ar[r]\&\mathcal G\ar[dr,dashed]\\
\&\&\&\&\Gamma
\end{tikzcd}}
\]
\margin{\uline{NB} L'hyp. $\bar k$ sép. clos dans $H^0(\mathcal X,\unsure{\mathcal O_{\mathcal X}})$ est \deleted{\unsure{automatiquement}} \inserted{stable} \unsure{par} passage à $\widetilde{\mathcal X},\mathcal X'$ et la \unsure{composée}
\diagram{\begin{tikzcd}[ampersand replacement=\&]
\mathcal G'\ar[r,no head]\ar[rr,bend right=25]\&\mathcal G\ar[r]\&\Gamma
\end{tikzcd}}
est l'hom. \unsure{déduit} de $\operatorname{Aut}_k(\mathcal X')\longrightarrow\Gamma=\operatorname{Aut}(\unsure{\bar k/k})$}

Dire que $\mathcal G'\to\mathcal G$ est surj revient à dire que $\mathcal E'=\mathcal E$, i.e. $N$ invariant \uline{dans} $\mathcal E$, i.e. $N$ \deleted{invariant} fixe sous l'action de $\mathcal G\deleted{\unsure{_k}}$ opérant sur l'ens. des sous-groupes invariants fermés de $\pi_1$.

\uline{NB} Si $\pi_1$ est \deleted{\unsure{top.}} de type fini, il admet un syst. fond. de voisinages de $1$ \unsure{formé} de ss-groupes (ouverts, i.e. fermés d'indice fini) invariants stables par \uline{tt} autom \inserted{continu}, et en particulier \deleted{\unsure{pour}} invariants dans $\mathcal E$.

\deleted{\unsure{Soit}}

\deleted{\textcircled{2.3} Une donnée de descente à $k$ de la \unsure{situation|relation} $\mathcal X'\mathbin{-}\mathcal X$ sur $k$}
